\section{Phân tích và lựa chọn phương pháp tiềm năng}
\label{sec:selection}

\subsection{Tiêu chí lựa chọn}

Việc lựa chọn được chi phối bởi bối cảnh triển khai (deployment) hơn là chỉ bởi
hiệu năng trên các bộ đánh giá chuẩn. Nền tảng phải chạy được tại các bệnh viện
có tài nguyên tính toán trải từ một máy chủ GPU chuyên dụng đến một máy trạm dùng
chung duy nhất, và phải trả kết quả đủ nhanh để tác động đến thứ tự đọc của danh
sách chờ đọc (worklist). Do đó các tiêu chí được áp dụng là:

\begin{enumerate}[label=C\arabic*.,leftmargin=1.2cm]
    \item \textbf{Hiệu năng chẩn đoán}, trong đó độ nhạy (sensitivity) theo từng
          phân nhóm được xem trọng hơn độ chính xác tổng hợp.
    \item \textbf{Độ trễ suy luận} (inference latency), hướng tới dưới hai phút
          cho mỗi ca chụp tính từ lúc tiếp nhận đến khi có kết quả, vì một tín
          hiệu phân loại ưu tiên (triage) đến sau khi bác sĩ chẩn đoán hình ảnh
          đã đọc xong ca chụp thì không mang lại lợi ích nào.
    \item \textbf{Yêu cầu phần cứng} (hardware footprint), bắt buộc phải có
          phương án suy giảm để thực thi chỉ trên CPU cho các cơ sở hạn chế tài
          nguyên.
    \item \textbf{Khả năng có dữ liệu huấn luyện}, ưu tiên các cách phát biểu bài
          toán được hỗ trợ bởi các tập dữ liệu công khai sẵn có.
    \item \textbf{Khả năng diễn giải} (interpretability), vì một kết luận không
          được giải thích sẽ không được tin tưởng hay sử dụng.
\end{enumerate}

\subsection{Phương pháp được lựa chọn}

\textbf{Kiến trúc hai giai đoạn} được lựa chọn: một bộ mã hoá (encoder) CNN 2D
áp dụng trên từng lát cắt (slice), tiếp theo là một mô hình chuỗi (sequence
model) hai chiều trên các vector nhúng (embedding) lát cắt thu được, xuất ra xác
suất phân nhóm ở cả mức lát cắt và mức ca chụp.

Lựa chọn này xuất phát từ các tiêu chí chứ không phải từ tính mới lạ. Nó đạt
hiệu năng cạnh tranh với các kết quả tốt nhất đã công bố (C1) trong khi vẫn khả
thi trên phần cứng khiêm tốn, vì thành phần tốn kém là một backbone 2D có chi phí
trên mỗi lát cắt nhỏ và dễ dàng xử lý theo lô (C2, C3). Nó huấn luyện trên nhãn
mức lát cắt, thứ mà tập dữ liệu RSNA cung cấp với số lượng lớn, còn giai đoạn mô
hình chuỗi không đòi hỏi gán nhãn (annotation) bổ sung (C4). Cấu trúc hai giai
đoạn là một lợi thế cho triển khai chứ không đơn thuần là bất tiện khi huấn
luyện: bộ mã hoá và mô hình chuỗi có thể tách ra trên các máy khác nhau, và tại
một cơ sở hạn chế tài nguyên, riêng bộ mã hoá cũng có thể đóng vai trò một bộ
phân loại theo lát cắt tuy suy giảm nhưng vẫn hoạt động được.

Việc dựng đầu vào tuân theo Mục~\ref{sec:windowing}: các giá trị lưu trữ được
chuyển sang đơn vị Hounsfield theo Công thức~\ref{eq:hu}, sau đó ba cửa sổ (não,
dưới màng cứng, xương) được ghép thành ba kênh đầu vào.

\subsection{Giai đoạn phân vùng bổ trợ}

Chỉ riêng phân loại không cho ra thể tích khối máu tụ (haematoma volume), vốn là
một đầu vào trực tiếp cho các quyết định điều trị. Do đó một giai đoạn phân vùng
(segmentation) dùng U-Net được đề xuất như một lượt xử lý thứ hai tuỳ chọn, chỉ
kích hoạt trên những ca chụp được bộ phân loại đánh dấu là dương tính. Việc đặt
điều kiện trên kết quả phân loại dương tính giúp giới hạn chi phí của nó, vì nó
chỉ chạy trên một phần nhỏ của thông lượng (throughput), và khu biệt các kiểu
thất bại của nó vào những ca mà dù sao con người cũng sẽ xem lại.

\subsection{Khả năng diễn giải}

Tiêu chí C5 được đáp ứng bằng hai cơ chế có mức độ chặt chẽ khác nhau. Các bản
đồ nổi bật (saliency map) dựa trên gradient (Grad-CAM và các phương pháp kế
thừa) phủ một bản đồ nhiệt lên lát cắt để chỉ ra các vùng đã chi phối dự đoán;
chúng rẻ và không đòi hỏi thay đổi kiến trúc, nhưng là một mô tả thô và đôi khi
không trung thực về lập luận của mô hình, nên cần được trình bày như một chỉ báo
chú ý chứ không phải như một đường phân định tổn thương. Ở nơi giai đoạn phân
vùng được chạy, mặt nạ của nó là một lời giải thích trực quan vững chắc hơn
nhiều, vì đó là một khẳng định không gian trực tiếp mà bác sĩ chẩn đoán hình ảnh
có thể kiểm chứng hoặc bác bỏ.

Xác suất đã hiệu chuẩn (calibration) cũng là một phần không thể thiếu của khả
năng diễn giải. Các mạng huấn luyện với hàm cross-entropy thường quá tự tin
(overconfident) một cách hệ thống, nên cần hiệu chuẩn hậu kiểm trước khi hiển thị
bất kỳ độ tin cậy dạng số nào cho bác sĩ lâm sàng; một con số ``94\%'' chưa hiệu
chuẩn mời gọi một mức độ tin tưởng mà mô hình chưa xứng đáng có được.

\subsection{Rủi ro và biện pháp giảm thiểu}

\begin{table}[H]
\centering
\small
\begin{tabularx}{\textwidth}{@{}X X@{}}
\toprule
\textbf{Rủi ro} & \textbf{Biện pháp giảm thiểu} \\
\midrule
Độ nhạy kém với SAH, phân nhóm khó nhất đối với cả người đọc lẫn hệ thống tự động & Báo cáo và ngưỡng riêng theo từng phân nhóm; ghi rõ trong giao diện rằng kết quả âm tính không loại trừ SAH \\
Dịch chuyển phân phối giữa các máy chụp và các cơ sở & Giám sát phân phối đầu vào; hiệu chuẩn lại ngưỡng theo từng cơ sở trong giai đoạn nghiệm thu \\
Xác suất quá tự tin & Hiệu chuẩn hậu kiểm; hiển thị khoảng tin cậy thay vì giá trị điểm \\
Thiên lệch tự động hoá khi đọc song song & Ban đầu chỉ triển khai ở chế độ phân loại ưu tiên; ẩn kết quả cho đến khi người đọc đã hình thành nhận định ban đầu \\
Mệt mỏi vì cảnh báo do dương tính giả ở tỷ lệ hiện mắc thấp & Ngưỡng được tinh chỉnh theo giá trị tiên đoán dương ở tỷ lệ hiện mắc tại chỗ, không chỉ theo độ nhạy \\
\bottomrule
\end{tabularx}
\caption{Các rủi ro chính của phương pháp được lựa chọn và biện pháp giảm thiểu.}
\label{tab:risks}
\end{table}
